\chapter{Procesy odnowy: zliczanie zdarzeń bez własności braku pamięci}
\label{ch:renewal}
\paragraph{Odsyłacze międzytomowe.} Modele intensywności i eksponowania czasu obserwacji znajdziesz w tomie IV, rozdziale~\ref{IV-ch:glm}.


\section{Dlaczego proces Poissona nie opisuje każdego strumienia zdarzeń?}
W procesie Poissona po dowolnej deterministycznej chwili obserwacji oczekiwanie na
kolejne zdarzenie ma nadal rozkład wykładniczy. To wynika z własności
braku pamięci. Awaria mechanizmu, którego prawdopodobieństwo
uszkodzenia rośnie z wiekiem, nie ma jednak na ogół tej własności.
Potrzebujemy więc modelu z niezależnymi, ale niekoniecznie
wykładniczymi czasami między kolejnymi zdarzeniami.
\emph{Proces odnowy} opisuje liczbę ukończonych cykli, w których
po każdym odnowieniu system zaczyna działać jak nowy.

\begin{definicja}[Proces odnowy]
Niech $(T_n)_{n\ge1}$ będzie ciągiem niezależnych zmiennych
losowych o wspólnym prawie $\nu$ na $(0,\infty)$.
Niech $F(t)=\Prob(T_1\le t)$ będzie jego dystrybuantą,
a $\overline F(t)=1-F(t)$ funkcją przeżycia.
Na wspólnym mierzalnym zbiorze zerowym możemy
zastąpić wszystkie odstępy wartością $1$, aby
$0<T_n<\infty$ zachodziło punktowo dla każdego $n$.
Nie zmienia to rozkładu całego ciągu.
Definiujemy $S_0=0$,
$S_n=\sum_{k=1}^nT_k$ oraz, dla $t\ge0$,
\[
N_t:=\sup\{n\ge0:S_n\le t\}\in\N_0\cup\{\infty\}.
\]
Proces $N=(N_t)_{t\ge0}$ nazywamy \emph{procesem odnowy}.
Liczby $S_n$ to chwile odnowień, a $T_n=S_n-S_{n-1}$
to czasy pomiędzy odnowieniami. Gdy $T_n\sim\mathrm{Exp}(\lambda)$,
odzyskujemy proces Poissona o intensywności $\lambda$
na mocy konstrukcji z rozdziału~\ref{ch:poisson-process}.
\end{definicja}

\begin{twierdzenie}[Brak eksplozji i podstawowe prawdopodobieństwa]
Dla procesu odnowy z poprzedniej definicji istnieje
mierzalny zbiór pełnej miary, na którym $N_t<\infty$
jednocześnie dla wszystkich $t\ge0$. Na tym zbiorze
trajektorie są niemalejące, prawostronnie ciągłe,
mają granice lewostronne i na każdym ograniczonym
przedziale wykonują skończenie wiele skoków wielkości $1$.
Dla $n\in\N_0$ zachodzi
\[
\Prob(N_t\ge n)=\Prob(S_n\le t),\qquad
\Prob(N_t=n)=\Prob(S_n\le t)-\Prob(S_{n+1}\le t).
\]
Dla $n\ge1$ prawem $S_n$ jest $n$-krotny splot miar
$\nu^{*n}$. Jego dystrybuantę oznaczamy przez
$F^{*n}(t)=\nu^{*n}([0,t])$; przyjmujemy też
$\nu^{*0}=\delta_0$.
\end{twierdzenie}
\begin{proof}
Ponieważ $\Prob(T_1>0)=1$, istnieje $\varepsilon>0$
takie, że $p:=\Prob(T_1\ge\varepsilon)>0$.
Zmienne $I_k=\chi_{\{T_k\ge\varepsilon\}}$ są niezależne
i mają wspólny rozkład Bernoulliego z parametrem $p$.
Z mocnego prawa wielkich liczb
$n^{-1}\sum_{k=1}^n I_k\to p>0$ prawie na pewno.
Ponieważ $S_n\ge\varepsilon\sum_{k=1}^nI_k$, wynika stąd
$S_n\to\infty$ prawie na pewno; dlatego na jednym
zbiorze pełnej miary $N_t$ jest skończone dla każdego
skończonego $t$. Na każdym ograniczonym przedziale
trajektoria jest funkcją schodkową z izolowanymi
punktami skoku $S_n$. Tożsamość
$N_t=\sum_{n\ge1}\chi_{\{S_n\le t\}}$ dowodzi
mierzalności $N_t$, a włączenie punktu $S_n$
do kolejnego przedziału stałości daje prawostronną
ciągłość. Ścisła dodatniość odstępów zapewnia
skoki wielkości $1$ i $N_0=0$.
Równość zdarzeń $\{N_t\ge n\}=\{S_n\le t\}$ jest bezpośrednia;
odejmując prawdopodobieństwa kolejnych zagnieżdżonych zdarzeń,
uzyskujemy drugi wzór. Rozkład sumy niezależnych zmiennych
jest splotem ich miar probabilistycznych.
\end{proof}

W dalszej części przyjmujemy wersję bez eksplozji
dla wszystkich trajektorii: na mierzalnym zbiorze
$\{S_n\not\to\infty\}$ miary zero zastępujemy
cały ciąg odstępów jedynkami. Zachowujemy wszystkie
jego prawa skończeniewymiarowe oraz niezależność.

\begin{uwaga}[Warunek $T_n>0$ jest istotny]
Gdyby dopuszczano $\Prob(T_1=0)>0$, lecz nadal
$\Prob(T_1>0)>0$, nie doszłoby do eksplozji, ale mogłyby
wystąpić wielokrotne skoki w tej samej chwili.
Założenie dodatniości usuwa ten problem i upraszcza zapis.
Jeżeli zaś $T_1=0$ prawie na pewno, wszystkie odnowienia
nastąpiłyby w chwili zerowej, a $N_0=\infty$.
\end{uwaga}

\begin{stwierdzenie}[Skończona średnia liczby odnowień]
Dla każdego procesu odnowy z $T_1>0$ prawie na pewno i każdego
$t<\infty$ zachodzi $\E N_t<\infty$, nawet gdy $\E T_1=\infty$.
\end{stwierdzenie}
\begin{proof}
Ustalmy $\theta>0$. Liczba
$r=\E e^{-\theta T_1}$ spełnia $0<r<1$.
Z nierówności Markowa i niezależności
\[
\Prob(S_n\le t)=\Prob(e^{-\theta S_n}\ge e^{-\theta t})
\le e^{\theta t}\E e^{-\theta S_n}=e^{\theta t}r^n.
\]
Z tożsamości $N_t=\sum_{n\ge1}\chi_{\{S_n\le t\}}$
i twierdzenia Tonellego otrzymujemy
$\E N_t=\sum_{n\ge1}\Prob(S_n\le t)
\le e^{\theta t}r/(1-r)<\infty$.
\end{proof}

\begin{center}
\begin{tikzpicture}[x=1cm,y=0.55cm,>=Latex]
\draw[->] (0,0)--(11,0) node[right] {$t$};
\draw[->] (0,0)--(0,5.8) node[above] {$N_t$};
\draw[very thick,blue!70!black]
(0,0)--(1.6,0) (1.6,1)--(3.8,1) (3.8,2)--(4.8,2)
(4.8,3)--(8,3) (8,4)--(10.2,4);
\foreach \x/\y in {1.6/0,3.8/1,4.8/2,8/3}{
 \draw[blue!70!black,very thick] (\x,\y)--(\x,{\y+1});
 \fill[blue!70!black] (\x,{\y+1}) circle (2pt);
 \draw[dashed,gray] (\x,0)--(\x,\y);
}
\foreach \x/\t in {1.6/1,3.8/2,4.8/3,8/4}{
 \node[below] at (\x,0) {$S_{\t}$};
}
\foreach \y in {1,2,3,4}{\node[left] at (0,\y) {$\y$};}
\end{tikzpicture}
\par\smallskip\small Schemat procesu odnowy. Odstępy
$T_1=S_1$, $T_2=S_2-S_1$, $T_3=S_3-S_2$ nie muszą mieć rozkładu
wykładniczego; dlatego liczba odnowień na rozłącznych przedziałach
zwykle nie jest niezależna.
\end{center}

\begin{przyklad}[Deterministyczny cykl]
Jeżeli $T_n\equiv c>0$, to $S_n=nc$ i
$N_t=\lfloor t/c\rfloor$. W modelu Poissona
$\Var(N_t)=\E N_t$, tutaj zaś $\Var(N_t)=0$.
Proces odnowy obejmuje więc zarówno modele losowe, jak i
całkowicie regularne cykle. Przykład pokazuje, że sam parametr
$1/\E T_1$ nie opisuje charakteru fluktuacji.
Ten deterministyczny proces jest jednak procesem Markowa
(w sensie niejednorodnym w czasie): znając aktualną chwilę,
znamy bez losowości wszystkie jego przyszłe wartości.
Nie może więc służyć jako dowód, że ogólny proces odnowy
nie ma własności Markowa. Rzeczywisty kontrprzykład podajemy dalej.
\end{przyklad}


\begin{twierdzenie}[Proces odnowy nie musi być procesem Markowa]
Niech $(T_n)_{n\ge1}$ będą niezależnymi zmiennymi, dla których
$\Prob(T_n=1)=\Prob(T_n=2)=1/2$, i niech
$N_t=\sup\{n\ge0:\sum_{k=1}^nT_k\le t\}$.
Wtedy proces $(N_t)_{t\ge0}$ \emph{nie jest} procesem Markowa
względem swojej filtracji naturalnej
$\mathcal F_t^N=\sigma(N_s:0\le s\le t)$.
\end{twierdzenie}
\begin{proof}
Spójrzmy na chwile $1$, $5/2$ i $13/4$.
Rozważmy dwa zdarzenia należące do $\mathcal F_{5/2}^N$:
\[
A=\{N_1=1,N_{5/2}=1\},\qquad
D=\{N_1=0,N_{5/2}=1\}.
\]
Oba mają dodatnie prawdopodobieństwo.
Na $A$ koniecznie $T_1=1$ i $T_2=2$: gdyby
$T_2=1$, druga odnowa nastąpiłaby już w chwili $2$.
Zatem $S_2=3<13/4$ i
$\Prob(N_{13/4}\ge2\mid A)=1$.
Na $D$ mamy $T_1=2$, a $T_2$ nadal przyjmuje wartości
$1$ i $2$ z jednakowymi prawdopodobieństwami,
niezależnie od $D$. Stąd
$\Prob(N_{13/4}\ge2\mid D)=1/2$.

Na obu zdarzeniach obecny stan $N_{5/2}$ jest taki sam,
mianowicie równy $1$. Gdyby $N$ miał własność Markowa,
warunkowe prawdopodobieństwo zdarzenia
$\{N_{13/4}\ge2\}$ przy całej informacji do chwili $5/2$
byłoby funkcją wyłącznie $N_{5/2}$. Jego średnie
warunkowe na $A$ i $D$ musiałyby więc być równe.
Obliczone wartości $1$ i $1/2$ przeczą temu warunkowi.
\end{proof}

\begin{uwaga}[Stan trzeba czasem poszerzyć]
Sama liczba odnowień $N_t$ nie mówi, ile czasu upłynęło
od ostatniej odnowy. Wprowadzamy wiek
$A_t=t-S_{N_t}$ i czas resztkowy $R_t=S_{N_t+1}-t$. To właśnie brakująca informacja w
powyższym dowodzie. Para $(N_t,A_t)$, gdzie
$A_t=t-S_{N_t}$, jest naturalnym kandydatem na stan
procesu Markowa: po znajomości wieku rozkład czasu
pozostałego do nowego cyklu można obliczyć z ogona
rozkładu odstępów. Własność Markowa takiej pary
wymaga osobnego sformułowania i dowodu; jej nie
utożsamiamy z własnością Markowa samego $N_t$.
\end{uwaga}


\begin{twierdzenie}[Wiek odnowy przywraca własność Markowa]
\label{thm:wiek-markow}
Niech $T_1,T_2,\ldots$ będą niezależnymi, jednakowo rozłożonymi
zmiennymi o wartościach w $(0,\infty)$, o dystrybuancie $F$,
a $N_t$ i $A_t=t-S_{N_t}$ oznaczają liczbę odnowień i wiek
bieżącego cyklu. Wówczas $Y_t=(N_t,A_t)$ jest jednorodnym
w czasie procesem Markowa względem filtracji naturalnej
$\calF_t^N=\sigma(N_s:0\le s\le t)$.
Dokładniej, warunkowe prawo czasu resztkowego
względem całej $\calF_t^N$ zależy tylko od $A_t$:
jego funkcja przeżycia dla $u\ge0$ wynosi
$\overline F(A_t+u)/\overline F(A_t)$ prawie na pewno,
a po następnym odnowieniu proces jest zbudowany z nowych,
niezależnych odstępów o prawie $F$. Opis ten nie zależy od $t$.
\end{twierdzenie}
\begin{proof}
Ustalmy deterministyczne $t$. Z trajektorii $(N_s)_{0\le s\le t}$
można odzyskać wszystkie chwile skoków $S_k\le t$, a więc
$N_t$ i $S_{N_t}$ są $\calF_t^N$-mierzalne. Na zdarzeniu
$\{N_t=n\}$ znamy $T_1,\ldots,T_n$, natomiast o $T_{n+1}$
wiemy tylko, że $T_{n+1}>t-S_n=A_t$. Istotnie,
$\{N_t=n\}=\{S_n\le t,\ T_{n+1}>t-S_n\}$.
Ściślej, na zdarzeniu $\{N_t=n\}$ ślad
$\calF_t^N$ jest równy śladowi
$\sigma(T_1,\ldots,T_n)$: każda obserwacja $N_s$,
$s\le t$, jest tam funkcją $S_1,\ldots,S_n$.
Odwrotnie, chwile $S_k$ dla $k\le n$ odzyskujemy
jako kresy dolne chwil osiągnięcia poziomu $k$,
sprawdzając chwile wymierne do $t$ oraz samą chwilę $t$.
Dla $n=0$ informacja na tym zdarzeniu jest trywialna.

Dla ograniczonej borelowskiej $\psi:(0,\infty)\to\R$
definiujemy borelowską funkcję
\[
g_\psi(a)=\frac{\int_{(a,\infty)}\psi(v-a)\,\mathrm d\nu(v)}
{\overline F(a)},\qquad \overline F(a)>0.
\]
Mierzalność licznika wynika z twierdzenia Fubiniego
dla parametru $a$. Na zdarzeniu $\{N_t=n\}$
warunkowe prawo $R_t$ spełnia
$\E[\psi(R_t)\mid\calF_t^N]=g_\psi(A_t)$
prawie na pewno. Aby zweryfikować tę równość,
dla ograniczonej borelowskiej funkcji $H$ wektora
$(T_1,\ldots,T_n)$ całkujemy po jego prawie:
\[
\E[H\chi_{\{N_t=n\}}\psi(R_t)]
=\E\left[H\chi_{\{S_n\le t\}}
\int_{(t-S_n,\infty)}\psi(v-(t-S_n))\,\mathrm d\nu(v)\right].
\]
Zastąpienie całki przez
$\overline F(t-S_n)g_\psi(t-S_n)$ daje
$\E[H\chi_{\{N_t=n\}}g_\psi(A_t)]$.
Równość śladów $\sigma$-ciał i sumowanie po $n$
uzasadniają warunkowanie względem całej historii.
Nie dzielimy przez prawdopodobieństwo zdarzenia
$\{A_t=a\}$.
Na zdarzeniu $\{N_t=n\}$ przyszłe odstępy
$T_{n+2},T_{n+3},\ldots$ zachowują prawo
produktowe $\nu^{\otimes\N}$ i są warunkowo niezależne
od $R_t$: powyższy rachunek powtarzamy z dodatkowym
czynnikiem będącym ograniczoną funkcją skończonej
liczby przyszłych odstępów. Niezależność pozwala
wyłączyć jej wartość oczekiwaną przed całkę;
argument $\pi$--$\lambda$ rozszerza wynik na cały ciąg. Po upływie resztkowego
czasu kolejna liczba odnowień oraz wiek procesu są zatem
jednoznacznymi funkcjami tych niezależnych odstępów i
aktualnej pary $(n,a)$. Dla dowolnej ograniczonej funkcji
borelowskiej $h$ na $\N_0\times[0,\infty)$ istnieje więc
funkcja $P_uh(n,a)$, która nie zależy od chwili $t$ i spełnia
$\E[h(Y_{t+u})\mid\calF_t^N]=P_uh(Y_t)$.
To jest żądana jednorodna własność Markowa.
Co istotne, nie wystarczy znać tylko rozkładu $R_t$:
po następnej odnowie całą przyszłość wyznaczają nowe,
niezależne odstępy. Właśnie ta niezależność pozwala
przejść od opisu pierwszego kolejnego skoku do dowolnej
liczby dalszych odnowień.
W szczególności, jeżeli do chwili $t+u$ nie ma nowej odnowy,
to $Y_{t+u}=(n,a+u)$; prawdopodobieństwo tego zdarzenia
jest równe $\overline F(a+u)/\overline F(a)$ przy
$\overline F(a)>0$. Zdarzenie $\{\overline F(A_t)=0\}$
ma jednak prawdopodobieństwo zero: dla każdego $n$ z niezależności
$T_{n+1}$ od $(T_1,\ldots,T_n)$ wynika
\[
\begin{aligned}
&\Prob\bigl(N_t=n,\,\overline F(A_t)=0\bigr)\\
&=\E\bigl[\chi_{\{S_n\le t,\,\overline F(t-S_n)=0\}}
     \Prob(T_{n+1}>t-S_n\mid T_1,\ldots,T_n)\bigr]=0.
\end{aligned}
\]
Sumując po $n$, uzasadniamy, że iloraz wystarcza na całym
zdarzeniu pełnej miary. Naturalną przestrzenią stanów jest
$E=\N_0\times\{a\ge0:\overline F(a)>0\}$.
Dla każdego jej punktu jądro pierwszego odstępu
jest dane powyższą całką, a dalsze odstępy mają
prawo produktowe $\nu^{\otimes\N}$.
Liczbę nowych odnowień i wiek w chwili $u$
otrzymujemy jako mierzalne funkcje tych odstępów
(przez zdarzenia wyznaczone sumami częściowymi).
Całkowanie względem tak skonstruowanego jądra daje
borelowską funkcję $P_uh$.
Po pierwszym skoku wiek zeruje się; jeśli skoku nie
ma, warunkowanie po przeżyciu dodatkowego czasu
$v$ daje ogon $\overline F(a+v+u)/\overline F(a+v)$.
Ten sam argument odnowy zastosowany do procesu
startującego w dowolnym $(n,a)\in E$ daje
$P_vP_u=P_{v+u}$.
\end{proof}

\begin{zadanie}[Dlaczego sam licznik nie wystarcza?]
Niech odstępy między odnowieniami mają rozkład jednostajny
na $[1,2]$. Ustalmy $t=2$ i rozważmy dwa możliwe przebiegi
z $N_t=1$: w pierwszym ostatnia odnowa nastąpiła w chwili
$t-1/4$, w drugim w chwili $t-3/4$. Porównać prawdopodobieństwa
braku odnowy w następnych $1/2$ jednostki czasu.

\emph{Rozwiązanie.} Dla wieku $a$ ogon pozostałego czasu
wynosi $\overline F(a+1/2)/\overline F(a)$.
Przy $a=1/4$ mamy $\overline F(3/4)=\overline F(1/4)=1$,
więc szansa braku odnowy jest $1$. Przy $a=3/4$ mamy
$\overline F(5/4)=3/4$ i $\overline F(3/4)=1$, zatem
szansa wynosi $3/4$. Zdarzenia warunkujące dokładnym
wiekiem mogą mieć prawdopodobieństwo zero; obliczenie
dotyczy regularnej wersji prawa warunkowego. Właśnie z tego
powodu para $(N_t,A_t)$ przenosi więcej informacji niż $N_t$.
\end{zadanie}

\section{Równanie odnowy i jego jednoznaczność}

Dla lokalnie ograniczonej funkcji borelowskiej $g$
na $[0,\infty)$ i lokalnie skończonej miary dodatniej
$V$ na tej półprostej piszemy
$(g*V)(t)=\int_{[0,t]}g(t-s)\,\mathrm dV(s)$.
Miara odnowy
\[
U=\sum_{n=0}^\infty \nu^{*n}
\]
obejmuje atom $\delta_0$ odpowiadający chwili początkowej.
Jest lokalnie skończona, gdyż $U([0,t])=1+\E N_t$.
Przeliczalna suma miar jest miarą z twierdzenia Tonellego.

\begin{twierdzenie}[Ogólne równanie odnowy]
\label{thm:renewal-equation}
Dla każdej lokalnie ograniczonej funkcji borelowskiej $g$
równanie
\[
u=g+u*\nu
\]
ma dokładnie jedno lokalnie ograniczone rozwiązanie
borelowskie, dane przez $u=g*U$.
\end{twierdzenie}
\begin{proof}
Dla $t\le T$ szereg
$\sum_{n\ge0}(g*\nu^{*n})(t)$ jest bezwzględnie
i jednostajnie zbieżny, gdyż suma oszacowań modułów
nie przekracza
$\sup_{[0,T]}|g|\,U([0,T])<\infty$.
Każdy składnik jest borelowski na mocy twierdzenia
Fubiniego. Z łączności splotu (Fubini na ograniczonym
trójkącie sumowania) i bezwzględnej zbieżności
otrzymujemy $(g*U)*\nu=g*(U-\delta_0)$,
co dowodzi istnienia.
Różnica $w$ dwóch rozwiązań spełnia
$w=w*\nu=w*\nu^{*n}$ dla każdego $n$.
Stąd $|w(t)|\le\sup_{[0,t]}|w|\,\Prob(S_n\le t)\to0$.
Dowolność $t$ daje jednoznaczność.
\end{proof}

\begin{twierdzenie}[Równanie dla średniej liczby odnowień]
Niech $N$ będzie procesem odnowy o dystrybuancie
odstępów $F$. Funkcja $m(t)=\E N_t$ jest skończona
dla $t\ge0$ i spełnia równanie
\[
m(t)=F(t)+\int_{(0,t]}m(t-s)\,\mathrm d\nu(s),\qquad t\ge0.
\]
Wśród lokalnie ograniczonych funkcji borelowskich
$u:[0,\infty)\to\R$ jest to jedyne rozwiązanie tego równania.
\end{twierdzenie}
\begin{proof}
Pierwsza odnowa następuje w chwili $T_1$. Jeżeli $T_1>t$,
liczba odnowień jest zero. Jeśli $T_1=s\le t$,
to pierwszą odnowę już zaliczamy, a kolejna liczba
odnowień ma rozkład $N'_{t-s}$, gdzie proces $N'$ jest
zbudowany z niezależnych odstępów $T_2,T_3,\ldots$
i ma to samo prawo co $N$. Stąd, warunkując względem
$T_1$, otrzymujemy
$m(t)=\E[\chi_{\{T_1\le t\}}(1+m((t-T_1)_+))]
=F(t)+\int_{(0,t]}m(t-s)\,\mathrm d\nu(s)$.
Dla ścisłości w drugim składniku używamy niezależności
$T_1$ od całego ciągu przyszłych odstępów, a nie tylko
równości rozkładów poszczególnych $T_n$.

Załóżmy, że $u_1,u_2$ są lokalnie ograniczonymi rozwiązaniami.
Dla $d=u_1-u_2$ mamy $d=d*\nu$,
a przez iterację $d=d*\nu^{*n}$.
Funkcja $m$ jest lokalnie ograniczona i borelowska,
ponieważ jest niemalejąca i skończona.
Można więc również zastosować
twierdzenie~\ref{thm:renewal-equation}.
Dla ustalonego $t$ niech $M_t=\sup_{0\le s\le t}|d(s)|<\infty$.
Wówczas
$|d(t)|\le M_t F^{*n}(t)=M_t\Prob(S_n\le t)\to0$,
bo $S_n\to\infty$ p.n. Stąd $d(t)=0$ i dowolność $t$
kończy dowód.
\end{proof}

\begin{przyklad}[Proces Poissona jako rozwiązanie równania odnowy]
Jeżeli $T_1\sim\mathrm{Exp}(\lambda)$, to $m(t)=\lambda t$.
Możemy sprawdzić ten wynik bez odwoływania się do rozkładu Poissona:
$F(t)=1-e^{-\lambda t}$, zatem prawa strona równania wynosi
\[
1-e^{-\lambda t}
+\int_0^t\lambda(t-s)\lambda e^{-\lambda s}\,\mathrm d\lambda_1(s)
=1-e^{-\lambda t}+\lambda t-(1-e^{-\lambda t})
=\lambda t.
\]
Z jednoznaczności rozwiązania równania odnowy otrzymujemy
$\E N_t=\lambda t$.
\end{przyklad}

\section{Mocne i elementarne twierdzenie odnowy}
\begin{twierdzenie}[Mocne prawo dla liczby odnowień]
Niech $N$ będzie procesem odnowy z $T_n>0$ p.n.
i $\mu=\E T_1\in(0,\infty]$.
Wtedy
\[
\frac{N_t}{t}\longrightarrow
\begin{cases}1/\mu,&\mu<\infty,\\0,&\mu=\infty,\end{cases}
\qquad\text{prawie na pewno, gdy }t\to\infty.
\]
\end{twierdzenie}
\begin{proof}
Najpierw przyjmijmy $0<\mu<\infty$.
Z MPWL $S_n/n\to\mu$ prawie na pewno.
Na zdarzeniu tej zbieżności $N_t\to\infty$ i
$S_{N_t}\le t<S_{N_t+1}$. Po podzieleniu nierówności
przez $N_t$ otrzymujemy
\[
\frac{S_{N_t}}{N_t}\le \frac{t}{N_t}
<\frac{N_t+1}{N_t}\frac{S_{N_t+1}}{N_t+1}.
\]
Obie skrajne wielkości dążą do $\mu$, więc
$t/N_t\to\mu$, a zatem $N_t/t\to1/\mu$.
Jeśli $\mu=\infty$, dla dowolnego $K>0$ zmienne
$T_n^{(K)}=\min(T_n,K)$ są niezależne, jednakowo rozłożone
i całkowalne. MPWL daje
$\liminf_{n\to\infty}S_n/n\ge\E T_1^{(K)}$.
Względem przeliczalnego ciągu $K\in\N$ i po przejściu
$K\to\infty$ twierdzenie o zbieżności monotonicznej
daje $S_n/n\to\infty$ prawie na pewno.
Ponownie $S_{N_t}\le t$ i $N_t\to\infty$, więc
$t/N_t\ge S_{N_t}/N_t\to\infty$; stąd $N_t/t\to0$.
\end{proof}

\begin{lemat}[Tożsamość Walda dla chwili przekroczenia]
Niech $T_1,T_2,\ldots$ będą niezależne i jednakowo rozłożone,
$T_n>0$ prawie na pewno, z $\mu=\E T_1<\infty$.
Dla $t\ge0$ połóżmy $K_t=N_t+1=\inf\{n\ge1:S_n>t\}$.
Wówczas $\E K_t<\infty$ i
$\E S_{K_t}=\mu\E K_t$.
\end{lemat}
\begin{proof}
Skończoność $\E K_t$ wynika z wcześniej udowodnionego
oszacowania $\E N_t<\infty$. Zauważmy, że
$\{K_t\ge i\}=\{S_{i-1}\le t\}$ zależy tylko od
$T_1,\ldots,T_{i-1}$, zatem jest niezależne od $T_i$.
Ponieważ wszystkie składniki są nieujemne, twierdzenie Tonellego
pozwala zapisać
\[
\E S_{K_t}=\E\sum_{i=1}^\infty T_i\chi_{\{K_t\ge i\}}
=\sum_{i=1}^\infty\mu\Prob(K_t\ge i)
=\mu\E K_t.
\]
Nie używamy tutaj błędnej zamiany $K_t$ na $N_t$:
liczba $N_t$ na ogół nie jest czasem zatrzymania
względem $\mathcal G_i=\sigma(T_1,\ldots,T_i)$,
gdyż $\{N_t\le i\}=\{S_{i+1}>t\}$ może zależeć
od niezaobserwowanego jeszcze $T_{i+1}$.
Na przykład dla odstępów $1$ lub $2$ z dodatnimi
prawdopodobieństwami i $t=3/2$, zdarzenie
$\{N_t=0\}=\{T_1=2\}$ nie należy do trywialnej
$\mathcal G_0$.
\end{proof}

\begin{twierdzenie}[Elementarne twierdzenie odnowy]
Przy założeniach definicji procesu odnowy niech
$\mu=\E T_1\in(0,\infty]$. Wtedy
\[
\frac{\E N_t}{t}\longrightarrow
\begin{cases}1/\mu,&\mu<\infty,\\0,&\mu=\infty,\end{cases}
\qquad\text{gdy }t\to\infty.
\]
\end{twierdzenie}
\begin{proof}
Dla $\mu<\infty$ niech $K_t=N_t+1$.
Ponieważ $S_{K_t}>t$, z lematu Walda dostajemy
$\mu(\E N_t+1)=\E S_{K_t}\ge t$,
czyli $\E N_t/t\ge1/\mu-1/t$.

Dla oszacowania od góry ustalmy $L>0$ i połóżmy
$T_i^{(L)}=\min(T_i,L)$, $\mu_L=\E T_1^{(L)}>0$.
Odnowienia w modelu z obciętymi odstępami występują
nie później, więc $N_t\le N_t^{(L)}$.
W tym modelu chwila pierwszego przekroczenia
$K_t^{(L)}=N_t^{(L)}+1$ spełnia
$t<S_{K_t^{(L)}}^{(L)}\le t+L$.
Tożsamość Walda (stosuje się do obciętych odstępów) daje
$\mu_L(\E N_t^{(L)}+1)\le t+L$. Wobec tego
\[
\limsup_{t\to\infty}\frac{\E N_t}{t}\le\frac{1}{\mu_L}.
\]
Przejście $L\to\infty$ i zbieżność monotoniczna
$\mu_L\uparrow\mu$ dają oszacowanie od góry $1/\mu$,
które wraz z oszacowaniem dolnym kończy dowód.
Jeśli $\mu=\infty$, stosujemy wyłącznie ostatnie
oszacowanie i otrzymujemy
$0\le\limsup_{t\to\infty}\E N_t/t\le1/\mu_L\to0$.
\end{proof}

\begin{uwaga}[Dlaczego silne prawo nie wystarcza?]
Zbieżność $N_t/t\to1/\mu$ prawie na pewno nie uprawnia
sama przez się do zamiany granicy i wartości oczekiwanej.
W ogólności potrzebna jest jednostajna całkowalność
albo odrębne oszacowania. W powyższym dowodzie rolę takich
oszacowań pełnią lemat Walda i obcięcie odstępów.
\end{uwaga}

\section{Twierdzenie Blackwella i kluczowe twierdzenie odnowy}

Elementarne twierdzenie odnowy opisuje średnią na
przedziale $[0,t]$, lecz nie mówi, co dzieje się
na przesuwanym przedziale o ustalonej długości.
Wprowadzamy rozróżnienie zależne od nośnika prawa
odstępu, a nie od tego, czy prawo ma gęstość.

\begin{definicja}[Rozkład arytmetyczny]
Prawo $\nu$ na $(0,\infty)$ nazywamy arytmetycznym,
jeżeli istnieje $h>0$ takie, że
$\nu(h\N)=1$. Największe takie $h$ nazywamy krokiem
arytmetycznym. W przeciwnym razie prawo jest
niearytmetyczne. Równoważnie, domknięta podgrupa
addytywna $\R$ generowana przez nośnik $\nu$
jest odpowiednio $h\Z$ lub $\R$.
\end{definicja}
Istotnie, każda domknięta niezerowa podgrupa
addytywna $\R$ ma tę postać: jeśli kres dolny
jej dodatnich elementów jest zerem, wielokrotności
dowolnie małych dodatnich elementów tworzą dowolnie
gęste siatki, więc domknięcie jest całą prostą.
Jeśli kres dolny jest dodatni, należy do podgrupy
przez domkniętość, a dzielenie z resztą dowodzi,
że wszystkie elementy są jego wielokrotnościami.

\begin{lemat}[Jednostajne ograniczenie miary odnowy]
\label{lem:renewal-local-bound}
Dla każdego $a>0$ istnieje skończona stała $C_a$
taka, że $U([x,x+a])\le C_a$ dla wszystkich $x\in\R$,
gdzie $U$ przedłużamy przez zero na $(-\infty,0)$.
\end{lemat}
\begin{proof}
Dla $x>0$ niech $J=\inf\{n\ge0:S_n\ge x\}$.
Jest to skończony czas zatrzymania względem
$\sigma(T_1,\ldots,T_n)$.
Jeśli $S_J\le x+a$, odnowień w $[x,x+a]$ jest
nie więcej niż $1+N'_a$, gdzie $N'$ jest procesem
zbudowanym z odstępów po indeksie $J$.
Warunkując na zdarzeniach $\{J=j\}$, które zależą
tylko od pierwszych $j$ odstępów, widzimy,
że przyszłe odstępy mają to samo prawo produktowe
co pierwotny ciąg. Zatem wartość oczekiwana
liczby odnowień jest nie większa niż $1+\E N_a$.
Dla $x\le0$ to samo ograniczenie wynika bezpośrednio
z $[x,x+a]\cap[0,\infty)\subseteq[0,a]$.
\end{proof}

\begin{twierdzenie}[Twierdzenie Blackwella]
\label{thm:blackwell-renewal}
Niech $\mu=\E T_1\in(0,\infty]$ i przyjmijmy
$1/\infty=0$.
Jeżeli $\nu$ jest niearytmetyczne, to dla każdego $a>0$
\[
\lim_{t\to\infty}U((t,t+a])=\frac a\mu.
\]
Jeżeli $\nu$ jest arytmetyczne o kroku $h$, to
\[
\lim_{n\to\infty}U(\{nh\})=\frac h\mu.
\]
\end{twierdzenie}
\begin{proof}
Podamy dowód metodą przesuwanych miar.
Z lematu~\ref{lem:renewal-local-bound} miary
$U_t(D)=U(t+D)$ mają jednostajnie ograniczoną
masę na każdym ograniczonym przedziale.
Z każdego ciągu $t_j\to\infty$ można wybrać
podciąg zbieżny mgliście (tj. w sensie zbieżności
całek wszystkich funkcji ciągłych o zwartym nośniku) do pewnej
lokalnie skończonej miary $V$.
Dla ścisłości wybieramy przeliczalny zbiór
funkcji gęsty w normie supremum w każdej
przestrzeni funkcji ciągłych o nośniku w $[-m,m]$,
$m\in\N$. Argument diagonalny daje zbieżność
całek na tym zbiorze, a lokalne ograniczenia
mas rozszerzają ją na całe $C_c(\R)$.
Dodatni funkcjonał graniczny reprezentuje
miara Radona na mocy twierdzenia Riesza--Markowa.
Granica $V$ również ma jednostajnie ograniczone
masy przedziałów jednostkowych: ograniczamy
ich funkcje charakterystyczne od góry przez
ciągłe funkcje o wartościach w $[0,1]$ i nośnikach
w przedziałach o długości $3$.

Z równania miar $U=\delta_0+U*\nu$ otrzymujemy
$V=V*\nu$. Rzeczywiście, dla $f\in C_c(\R)$
\[
\int f\,\mathrm dU_{t_j}
=f(-t_j)+\int_{(0,\infty)}
\left(\int f(z+s)\,\mathrm dU_{t_j}(z)\right)\mathrm d\nu(s).
\]
Wewnętrzne całki są jednostajnie ograniczone
niezależnie od $s,j$ dzięki lematowi
\ref{lem:renewal-local-bound}. Zbieżność mglista
dla ustalonego $s$ i zbieżność zmajoryzowana
dają równanie dla $V$.

Ustalmy $f\in C_c(\R)$ i połóżmy
$H(x)=\int f(x-y)\,\mathrm dV(y)$.
Jest to ograniczona ciągła funkcja
(na mocy jednostajnej ciągłości $f$ i
jednostajnych ograniczeń mas).
Z $V=V*\nu$ wynika
$H(x)=\int H(x-s)\,\mathrm d\nu(s)$.
Dla ustalonego $x$ ciąg $H(x-S_n)$ jest
ograniczonym martyngałem względem odstępów.
Ma więc granicę $L_x$ prawie na pewno i w $L^1$
(z twierdzenia o zbieżności martyngałów ograniczonych).
Granica jest niezmiennicza względem skończonych
permutacji odstępów, ponieważ od pewnego indeksu
takie permutacje nie zmieniają sum częściowych.

Przypomnijmy krótki argument zero-jedynkowy.
Jeżeli zdarzenie $A$ zależne od niezależnych,
jednakowo rozłożonych zmiennych jest niezmiennicze
względem skończonych permutacji, to dla każdego
$\varepsilon>0$ można przybliżyć je zdarzeniem $D$
zależnym od pierwszych $m$ zmiennych tak, że
$\Prob(A\mathbin{\triangle}D)<\varepsilon$.
Istnienie takiego przybliżenia wynika z faktu,
że zdarzenia przybliżalne w mierze tworzą
$\sigma$-ciało zawierające cylindry (dla sumy
przeliczalnej najpierw bierzemy sumę skończoną).
Zamieniając pierwsze $m$ zmiennych z następnymi
$m$, dostajemy niezależną kopię $D'$ zdarzenia $D$
i nadal $\Prob(A\mathbin{\triangle}D')<\varepsilon$.
Stąd $|\Prob(A)-\Prob(D\cap D')|<2\varepsilon$,
a $\Prob(D\cap D')=\Prob(D)^2$.
Przejście $\varepsilon\downarrow0$ daje
$\Prob(A)=\Prob(A)^2$.
Stosując ten argument do zdarzeń
$\{L_x\le q\}$ dla wymiernych $q$, widzimy,
że $L_x$ jest stała p.n.
Możemy pracować na kanonicznej przestrzeni
produktowej odstępów i zdefiniować granicę
przez $\limsup$, dzięki czemu niezmienniczość
jest punktowa. Ponieważ $\E L_x=H(x)$,
z własności martyngału i zbieżności w $L^1$
wynika $H(x-T_1)=\E[L_x\mid T_1]=H(x)$ p.n.
Ciągłość $H$ daje $H(x-s)=H(x)$ dla każdego
$s$ z nośnika $\nu$, dla każdego ustalonego $x$.
W konsekwencji każdy element domkniętej
podgrupy generowanej przez nośnik jest okresem $H$.
Ponieważ dotyczy to każdej $f\in C_c(\R)$,
miara $V$ jest niezmiennicza względem przesunięć
o elementy tej podgrupy.

W przypadku niearytmetycznym oznacza to
$V=c\lambda_1$ dla $c\ge0$. Istotnie,
niezmienniczość wymusza brak atomów
(lokalna skończoność), a masa przedziału
jest addytywną niemalejącą funkcją jego długości,
zatem jest proporcjonalna do długości.
W przypadku arytmetycznym rozważamy wyłącznie
$t_j=n_jh$. Miary $U_{t_j}$ i ich granica
są skupione na domkniętym zbiorze $h\Z$,
a niezmienniczość względem $h$ daje
$V=c\sum_{k\in\Z}\delta_{kh}$.

Pozostaje ustalić $c$. Dokładnie jeden cykl
zawiera chwilę $t$, więc niezależność
$T_{n+1}$ od $S_n$ daje tożsamość
\[
1=\sum_{n\ge0}\Prob(S_n\le t<S_{n+1})
=\int_{[0,t]}\overline F(t-s)\,\mathrm dU(s).
\]
W zmiennej przesuniętej $z=s-t$ jest to
całka $\overline F(-z)$ po $[-t,0]$
względem $U_t$. Dla $\mu<\infty$
ogon tej całki po $z<-M$ jest jednostajnie
mały, ponieważ z lokalnego ograniczenia
jest nie większy niż
$C_1\sum_{k\ge\lfloor M\rfloor}\overline F(k)$,
a szereg jest zbieżny z
$\int_0^\infty\overline F(v)\,\mathrm d\lambda_1(v)=\mu$
(tożsamość wynika z twierdzenia Tonellego zastosowanego
do $T_1=\int_0^\infty\chi_{\{v<T_1\}}\,\mathrm d\lambda_1(v)$).
W przypadku niearytmetycznym na każdym
$[-M,0]$ możemy przejść do granicy:
funkcja monotoniczna $\overline F(-z)$ ma
co najwyżej przeliczalnie wiele nieciągłości,
które mają miarę $V$ zero. Uzasadnienie
dają górne i dolne sumy schodkowe na podziałach
o końcach bez atomów $V$ oraz zbieżność
mas tych przedziałów wynikająca z mglistej
zbieżności przez aproksymację ciągłą.
Otrzymujemy $1=c\mu$.
Gdy $\mu=\infty$, używamy wyłącznie obciętych
całek: $1\ge c\int_0^M\overline F(v)\,\mathrm d\lambda_1(v)$,
więc $c=0$.

W przypadku arytmetycznym na każdym skończonym
zbiorze punktów siatki masy poszczególnych punktów
zbiegają do $c$:
każdy punkt można izolować funkcją ciągłą
o nośniku w przedziale długości mniejszej niż $h$.
Analogicznie dla $\mu<\infty$ dostajemy
$1=c\sum_{k\ge0}\overline F(kh)=c\mu/h$,
gdzie ostatnia równość wynika z wzoru na średnią
dodatniej całkowitoliczbowej zmiennej $T_1/h$.
Przy nieskończonej średniej skończone sumy
wymuszają $c=0$.

Każda granica podciągowa jest więc taka sama.
W przypadku niearytmetycznym końce $(0,a]$
nie mają masy granicznej, więc zbieżność mglista
daje $U_{t_j}((0,a])\to a/\mu$.
W przypadku arytmetycznym izolujemy punkt zero
i dostajemy $U_{n_jh}(\{0\})\to h/\mu$.
Argument podciągowy kończy dowód.
\end{proof}

\begin{definicja}[Bezpośrednia całkowalność Riemanna]
Nieujemną funkcję $g:[0,\infty)\to\R$ nazywamy
bezpośrednio całkowalną w sensie Riemanna, jeżeli
dla przedziałów $I_{k,\delta}=[k\delta,(k+1)\delta)$
górne sumy
$\delta\sum_{k\ge0}\sup_{I_{k,\delta}}g$
są skończone dla wszystkich dostatecznie małych
$\delta>0$, a różnica górnej i dolnej sumy
dąży do zera przy $\delta\downarrow0$.
Dla funkcji ze znakiem wymagamy tej własności
od części dodatniej i ujemnej.
\end{definicja}

\begin{twierdzenie}[Kluczowe twierdzenie odnowy]
\label{thm:key-renewal}
Jeżeli $\nu$ jest niearytmetyczne, a $g$ jest
borelowska i bezpośrednio całkowalna w sensie
Riemanna, to
\[
\lim_{t\to\infty}(g*U)(t)
=\frac1\mu\int_0^\infty g(v)\,\mathrm d\lambda_1(v).
\]
W przypadku arytmetycznym o kroku $h$,
jeśli $\sum_{k\ge0}|g(kh)|<\infty$, to
\[
\lim_{n\to\infty}(g*U)(nh)
=\frac h\mu\sum_{k\ge0}g(kh).
\]
\end{twierdzenie}
\begin{proof}
Dla $g\ge0$ w przypadku niearytmetycznym
ustalmy $\delta>0$ z definicji.
Rozbijamy całkę względem $U_t$ na przedziały
$(- (k+1)\delta,-k\delta]$ i ograniczamy ją
górnymi oraz dolnymi wartościami $g$ na
$I_{k,\delta}$. Miara każdego takiego przedziału
dąży do $\delta/\mu$ na mocy twierdzenia
\ref{thm:blackwell-renewal}, a jest jednostajnie
ograniczona przez $C_\delta$ z lematu
\ref{lem:renewal-local-bound}. Przedłużenie
$U_t$ przez zero poza jego nośnikiem pozwala
sumować po wszystkich $k\ge0$.
Sumowalność górnych wartości daje zbieżność
sum przez zmajoryzowanie szeregu.
Dolna i górna granica splotu leżą więc między
dolną i górną sumą Riemanna podzielonymi przez
$\mu$. Przejście $\delta\downarrow0$ daje
wskazaną całkę; gdy $\mu=\infty$, już dla
jednego $\delta$ górna granica wynosi zero.
Dla funkcji ze znakiem stosujemy dowód do obu części.
W przypadku arytmetycznym splot w chwili $nh$
jest sumą $\sum_{k=0}^n g(kh)U(\{(n-k)h\})$.
Każdy czynnik miary zbiega do $h/\mu$,
a lokalne ograniczenie i bezwzględna sumowalność
pozwalają przejść do granicy w szeregu.
\end{proof}

\begin{przyklad}[Dlaczego arytmetyczność ma znaczenie?]
Dla $T_n=h$ punktowo miara odnowy jest
$\sum_{k\ge0}\delta_{kh}$.
Jeśli $0<a<h$, to $U((t,t+a])$ przyjmuje
wartości $0$ i $1$ w dowolnie późnych chwilach,
więc nie może zbiegać do $a/h$.
Natomiast $U(\{nh\})=1=h/\mu$.
Prawo o dwóch atomach $1,\sqrt2$ jest
niearytmetyczne, mimo że nie ma gęstości:
żadna dodatnia liczba nie jest wspólnym krokiem
dla obu tych liczb, ponieważ ich stosunek
jest niewymierny.
\end{przyklad}

\section{Wiek odnowy i czas pozostały do następnego zdarzenia}
\begin{definicja}[Wiek i czas resztkowy]
Dla $t\ge0$ oraz procesu odnowy o chwilach $S_n$ definiujemy
$A_t:=t-S_{N_t}$ oraz $R_t:=S_{N_t+1}-t$.
Nazywamy je odpowiednio \emph{wiekiem bieżącego cyklu}
i \emph{pozostałym czasem do następnego odnowienia}.
\end{definicja}

\begin{stwierdzenie}[Warunkowy rozkład czasu resztkowego]
Niech $\nu$ będzie prawem ściśle dodatnich odstępów,
a $F$ jego dystrybuantą.
Ustalmy $t\ge0$, $n\ge0$ oraz $s\in[0,t]$ takie, że
$\overline F(t-s)>0$. Wtedy możemy wybrać regularną wersję rozkładu
warunkowego względem $(N_t,S_n)$, która dla
warunków $N_t=n$, $S_n=s$ spełnia dla $u\ge0$
\[
\Prob(R_t>u\mid N_t=n,\ S_n=s)
=\frac{\overline F(t-s+u)}{\overline F(t-s)}.
\]
\end{stwierdzenie}
\begin{proof}
Na zdarzeniu $\{S_n\le t\}$ warunek $N_t=n$
jest równoważny $T_{n+1}>t-S_n$.
Zmienna $T_{n+1}$ jest niezależna od $S_n$, gdyż
$S_n$ zależy tylko od poprzednich odstępów.
Warunkujemy zatem zmienną $T_{n+1}$ na przekroczenie
progu $t-s$ i dostajemy iloraz
$\Prob(T_{n+1}>t-s+u)/\Prob(T_{n+1}>t-s)$.
Formalnie dla dowolnego zbioru borelowskiego $D$
\[
\Prob(N_t=n,S_n\in D,R_t>u)
=\int_{D\cap[0,t]}\overline F(t-s+u)\,\mathrm d\nu^{*n}(s).
\]
Z kolei miara rozkładu $S_n$ ograniczona
do zdarzenia $\{N_t=n\}$ wynosi
$\chi_{[0,t]}(s)\overline F(t-s)\,\mathrm d\nu^{*n}(s)$.
Całkowanie podanego ilorazu względem tej miary
daje dokładnie poprzednią równość, także dla $n=0$.
Jest to uzasadnienie regularnej wersji warunkowej
z twierdzenia~\ref{thm:wiek-markow}.
\end{proof}

\begin{przyklad}[Dwie zupełnie różne pamięci]
Dla odstępów wykładniczych $\overline F(v)=e^{-\lambda v}$,
więc iloraz wynosi $e^{-\lambda u}$, niezależnie od wieku
$t-s$. Stąd proces Poissona zapomina czas ostatniego zgłoszenia.
Dla deterministycznych odstępów $T_i\equiv c$ pozostały czas
jest natomiast równy $c-A_t$ we wszystkich chwilach,
w tym w chwilach odnowień, gdy $A_t=0$ i $R_t=c$. Znając wiek, dokładnie znamy czas kolejnej
odnowy. Nie należy więc oczekiwać własności Markowa dla
samego procesu $N_t$ w dowolnym procesie odnowy; często
naturalnym stanem jest para $(N_t,A_t)$.
\end{przyklad}

\section{Graniczny wiek, czas resztkowy i paradoks inspekcji}

\begin{twierdzenie}[Graniczny rozkład wieku i czasu resztkowego]
\label{thm:renewal-age-limit}
Załóżmy, że $\mu=\E T_1<\infty$ i prawo $\nu$
jest niearytmetyczne. Wówczas dla $a,b\ge0$
\[
\lim_{t\to\infty}\Prob(A_t\ge a,\ R_t>b)
=\frac1\mu\int_{a+b}^\infty\overline F(v)\,
\mathrm d\lambda_1(v).
\]
Para $(A_t,R_t)$ zbiega według rozkładu do pary
$(A_*,R_*)$, której obie współrzędne mają
rozkład równowagowy o gęstości
$\overline F(v)/\mu$ na $[0,\infty)$.
\end{twierdzenie}
\begin{proof}
Dla ustalonego $t\ge0$ dokładnie jeden indeks $n$
spełnia $S_n\le t<S_{n+1}$.
Na tym zdarzeniu warunki $A_t\ge a$, $R_t>b$
są równoważne $S_n\le t-a$ i $T_{n+1}>t+b-S_n$.
Z niezależności odstępów, sumowania po $n$ i
definicji $U$ otrzymujemy
\[
\Prob(A_t\ge a,R_t>b)
=\int_{[0,t]}\chi_{[a,\infty)}(t-s)
\overline F(t-s+b)\,\mathrm dU(s).
\]
Funkcja $g(v)=\chi_{[a,\infty)}(v)\overline F(v+b)$
jest bezpośrednio całkowalna w sensie Riemanna:
poza jednym punktem skoku jest niemalejąca
na początkowym przedziale (gdzie jest równa zeru)
i nierosnąca po $a$. Dla każdej siatki górna
suma jest skończona z całkowalności ogona,
a różnica górnej i dolnej sumy jest ograniczona
przez $2\delta$, z oszacowania całkowitej wariacji.
Twierdzenie~\ref{thm:key-renewal} daje żądaną granicę.
Aby zidentyfikować prawo graniczne, wybierzmy
$L$ o prawie
\[
\nu_{\rm len}(D)=\frac1\mu\int_D v\,\mathrm d\nu(v),
\]
oraz niezależną zmienną $Q$ jednostajną na $(0,1)$.
Połóżmy $A_*=QL$ i $R_*=(1-Q)L$.
Warunkując względem $L$ otrzymujemy
\[
\Prob(A_*\ge a,R_*>b)
=\frac1\mu\int_{(a+b,\infty)}(v-a-b)\,\mathrm d\nu(v)
=\frac1\mu\int_{a+b}^\infty\overline F(v)\,\mathrm d\lambda_1(v).
\]
Ostatnia równość wynika z twierdzenia Tonellego.
Dla $a=0$ lub $b=0$ dostajemy rozkłady brzegowe
o ciągłych dystrybuantach. Granice wspólnych
ogonów i ogonów brzegowych wyznaczają granice
dystrybuant łącznych (dla ścisłych nierówności
dotyczących $A_t$ stosujemy oszacowanie przez
progi $a$ i $a+\varepsilon$ oraz ciągłość
rozkładu granicznego). Stąd zbieżność według
rozkładu całej pary.
\end{proof}

\begin{wniosek}[Długość cyklu widzianego w późnej chwili]
Przy założeniach twierdzenia~\ref{thm:renewal-age-limit}
długość bieżącego cyklu $L_t=A_t+R_t=T_{N_t+1}$
zbiega według rozkładu do $L$ o prawie
$\nu_{\rm len}(\mathrm dv)=v\,\nu(\mathrm dv)/\mu$.
Ponadto
\[
\E L=\frac{\E T_1^2}{\mu},\qquad
\E A_*=\E R_*=\frac{\E T_1^2}{2\mu},
\]
z dopuszczeniem nieskończonych wartości.
\end{wniosek}
\begin{proof}
Zbieżność wynika z twierdzenia o odwzorowaniu
ciągłym zastosowanego do sumy współrzędnych.
Wzór na $\E L$ wynika z definicji jego prawa,
a niezależność $Q$ i $L$ oraz $\E Q=1/2$
dają pozostałe wzory przez Tonellego.
Nie twierdzimy na tej podstawie, że wartości
oczekiwane $L_t,A_t,R_t$ zbiegają: do tego
potrzebne jest dodatkowe uzasadnienie.
\end{proof}

\begin{uwaga}[Paradoks inspekcji]
Losowy cykl obserwowany w późnej deterministycznej
chwili nie ma na ogół prawa zwykłego odstępu $\nu$.
Cykl długości $v$ ma wagę proporcjonalną do $v$.
Jeżeli drugi moment odstępu jest skończony, to
\[
\E L=\mu+\frac{\Var(T_1)}{\mu}\ge\mu.
\]
Równość zachodzi jedynie dla odstępów stałych.
Interpretacja przez rozkład długościowo obciążony
obowiązuje również przy nieskończonym drugim
momencie, kiedy średnia obserwowanej długości
jest nieskończona.
\end{uwaga}

\begin{przyklad}[Odstępy jednostajne]
Jeżeli $T_1$ ma rozkład jednostajny na $[1,2]$,
to $\mu=3/2$ i $\E T_1^2=7/3$.
Graniczna długość obserwowanego cyklu ma gęstość
$2v/3$ na $[1,2]$, a jej średnia wynosi $14/9$.
Średni graniczny czas resztkowy wynosi $7/9$,
a nie $\mu/2=3/4$. Jego gęstość jest równa
$2/3$ na $[0,1)$ i $2(2-v)/3$ na $[1,2]$,
a poza tymi przedziałami zero.
\end{przyklad}

\begin{stwierdzenie}[Graniczny rozkład w przypadku arytmetycznym]
Niech $\nu$ będzie arytmetyczne o kroku $h$,
$\mu<\infty$. Ustalmy fazę $\theta\in[0,h)$.
Dla $a,b\ge0$ zachodzi
\[
\lim_{n\to\infty}\Prob(A_{nh+\theta}\ge a,R_{nh+\theta}>b)
=\frac h\mu\sum_{k\ge0}
\chi_{[a,\infty)}(kh+\theta)\overline F(kh+\theta+b).
\]
\end{stwierdzenie}
\begin{proof}
Wykorzystujemy dokładną tożsamość splotową
z dowodu twierdzenia~\ref{thm:renewal-age-limit}.
Wszystkie odnowienia są na siatce $h\N_0$,
więc dla $t=nh+\theta$ otrzymujemy sumę
\[
\sum_{k=0}^n
\chi_{[a,\infty)}(kh+\theta)\overline F(kh+\theta+b)
U(\{(n-k)h\}).
\]
Wagi ogonowe są sumowalne, gdyż
$\sum_{k\ge0}\overline F(kh)=\mu/h$.
Z twierdzenia~\ref{thm:blackwell-renewal}
i jednostajnego ograniczenia mas punktów siatki
wynika teza przez zmajoryzowanie szeregu.
W szczególności nie wolno w ogólności oczekiwać
jednej granicy dla wszystkich rzeczywistych $t$:
dla $T_1=h$ mamy dokładnie $A_{nh+\theta}=\theta$
i $R_{nh+\theta}=h-\theta$.
\end{proof}

\section{Proces odnowy w równowadze}

\begin{definicja}[Opóźniony proces odnowy]
W opóźnionym procesie odnowy pierwsze oczekiwanie
$D>0$ może mieć inne prawo niż następne odstępy.
Zakładamy niezależność $D$ od niezależnych odstępów
$T_1,T_2,\ldots$ o wspólnym prawie $\nu$.
Chwile odnowień wynoszą
$D+S_n$, $n\ge0$, gdzie $S_0=0$.
Gdy $\mu<\infty$, proces nazywamy równowagowym,
jeśli $D$ ma gęstość
$\overline F(v)/\mu$ dla $v\ge0$.
\end{definicja}

\begin{twierdzenie}[Stacjonarność przyrostów procesu równowagowego]
\label{thm:renewal-equilibrium}
Niech $\widetilde N_t$ zlicza odnowienia procesu
równowagowego na $(0,t]$.
Wówczas dla każdego $t,u\ge0$
\[
\E\widetilde N_t=\frac t\mu,\qquad
\Prob(\widetilde R_t>u)=
\frac1\mu\int_u^\infty\overline F(v)\,\mathrm d\lambda_1(v),
\]
gdzie $\widetilde R_t$ jest czasem do najbliższego
odnowienia ściśle po $t$. Ponadto proces
$(\widetilde N_{t+u}-\widetilde N_t)_{u\ge0}$
ma to samo prawo co $(\widetilde N_u)_{u\ge0}$.
Nie zakładamy niearytmetyczności $\nu$.
\end{twierdzenie}
\begin{proof}
Niech $\eta$ będzie prawem $D$. Miara zliczająca
oczekiwane odnowienia opóźnione wynosi
$\widetilde U=\eta*U$, przez niezależność i
sumowanie praw chwil $D+S_n$.
Ma ona gęstość względem miary Lebesgue'a:
z twierdzenia Tonellego
\[
\frac{\mathrm d\widetilde U}{\mathrm d\lambda_1}(t)
=\frac1\mu\int_{[0,t]}\overline F(t-s)\,\mathrm dU(s)
=\frac1\mu.
\]
Ostatnia równość to tożsamość wykorzystana
w dowodzie twierdzenia~\ref{thm:blackwell-renewal};
zachodzi dla każdego $t\ge0$.
Zatem $\E\widetilde N_t=\widetilde U((0,t])=t/\mu$.

Zdarzenie $\{\widetilde R_t>u\}$ zachodzi albo
gdy $D>t+u$, albo gdy jedno z odnowień
$D+S_n\le t$ jest ostatnim przed $t$ i następujący
po nim odstęp przekracza $t+u-D-S_n$.
Zdarzenia te są rozłączne. Z niezależności i
otrzymanej postaci $\widetilde U$ wynika
\[
\begin{aligned}
\Prob(\widetilde R_t>u)
&=\Prob(D>t+u)+
\int_{[0,t]}\overline F(t+u-s)\,\mathrm d\widetilde U(s)\\
&=\frac1\mu\left(
\int_{t+u}^\infty\overline F(v)\,\mathrm d\lambda_1(v)
+\int_u^{t+u}\overline F(v)\,\mathrm d\lambda_1(v)\right).
\end{aligned}
\]
To jest ogon prawa $D$.

Odstępy po pierwszym odnowieniu ściśle po $t$
są niezależne od $\widetilde R_t$ i mają prawo
produktowe $\nu^{\otimes\N}$.
Weryfikujemy to, rozbijając przestrzeń na
$\{D>t\}$ oraz $\{D+S_n\le t<D+S_{n+1}\}$,
$n\ge0$. Na każdym z tych zdarzeń czas resztkowy
i samo zdarzenie zależą tylko od odstępów do
najbliższego odnowienia włącznie; dalszy ciąg
jest niezależny od tych zmiennych i ma ustalone
prawo produktowe. Faktoryzacja dla ograniczonych
funkcji czasu resztkowego i cylindrów dalszych
odstępów, a następnie sumowanie oraz argument
$\pi$--$\lambda$, dowodzą niezależności.
Zatem proces liczby odnowień po chwili $t$
ma tę samą konstrukcję co pierwotny proces
równowagowy. Otrzymujemy równość praw procesów,
w szczególności wszystkich rozkładów
skończeniewymiarowych.
\end{proof}

\begin{uwaga}[Stacjonarne przyrosty a stacjonarność licznika]
Licznik $\widetilde N_t$ nie jest procesem
stacjonarnym: zaczyna się w zerze, a jego średnia
rośnie jak $t/\mu$. Stacjonarne są jego przyrosty,
czyli liczby odnowień na przesuwanych przedziałach
oraz ich wspólne rozkłady. Nie oznacza to
niezależności przyrostów. Dla stałych odstępów $h$
pierwsze oczekiwanie w równowadze jest jednostajne
na $(0,h)$, co losuje fazę siatki odnowień.
\end{uwaga}

\section{Nagrody w cyklach: twierdzenie odnowy z nagrodą}
\begin{twierdzenie}[Długookresowa średnia nagrody]
Niech $(T_n,R_n)_{n\ge1}$ będą niezależnymi parami o jednakowym
rozkładzie w $(0,\infty)\times\R$, gdzie
$0<\mu=\E T_1<\infty$ oraz $\E|R_1|<\infty$.
Nie wymagamy niezależności $T_n$ i $R_n$ wewnątrz pary.
Niech $S_n=\sum_{k=1}^nT_k$, $N_t=\max\{n:S_n\le t\}$ i
$C_t=\sum_{k=1}^{N_t}R_k$. Wówczas
\[
\frac{C_t}{t}\longrightarrow\frac{\E R_1}{\E T_1}
\qquad\text{prawie na pewno, gdy }t\to\infty.
\]
\end{twierdzenie}
\begin{proof}
Na wspólnym zdarzeniu pełnej miary mocne prawo wielkich liczb
dla obu współrzędnych par daje
$S_n/n\to\mu$ oraz $n^{-1}\sum_{k=1}^nR_k\to\rho:=\E R_1$.
Z twierdzenia odnowy $N_t\to\infty$ oraz
$N_t/t\to1/\mu$ prawie na pewno. Wobec tego
\[
\frac{C_t}{t}=
\left(\frac{1}{N_t}\sum_{k=1}^{N_t}R_k\right)
\frac{N_t}{t}\longrightarrow\rho/\mu.
\]
Jeżeli przez pewien początkowy czas $N_t=0$, przyjmujemy
$C_t=0$; nie ma to wpływu na granicę.
\end{proof}

\begin{zadanie}[Udział czasu naprawy maszyny]
Maszyna pracuje przez czas $U_n\sim\mathrm{Exp}(\alpha)$,
a następnie jest naprawiana przez czas
$V_n\sim\mathrm{Exp}(\beta)$. Załóżmy, że wszystkie
zmienne $(U_n,V_n)$ są niezależne pomiędzy cyklami,
a w każdym cyklu czasy pracy i naprawy są niezależne.
Obliczyć długookresowy udział czasu napraw.

\emph{Rozwiązanie.} Czas jednego cyklu wynosi
$T_n=U_n+V_n$, a nagrodą jest długość naprawy $R_n=V_n$.
Zatem $\E T_1=1/\alpha+1/\beta$ i $\E R_1=1/\beta$.
Twierdzenie daje udział napraw w ukończonych cyklach
\[
\frac{\E V_1}{\E(U_1+V_1)}
=\frac{1/\beta}{1/\alpha+1/\beta}
=\frac{\alpha}{\alpha+\beta}.
\]
Aby uzyskać udział rzeczywistego czasu napraw na przedziale
$[0,t]$, należy uwzględnić nieukończony cykl. Różnica
między nagrodą z ukończonych cykli a rzeczywistym czasem
napraw nie przekracza $T_{N_t+1}$. Ponieważ
$S_n/n\to\mu$, mamy
$T_n/n=(S_n/n)-((n-1)/n)(S_{n-1}/(n-1))\to0$ p.n.,
a stąd $T_{N_t+1}/t\to0$ po połączeniu z $N_t/t\to1/\mu$.
Dostajemy ten sam udział dla całego czasu $[0,t]$.
\end{zadanie}


\section{Transformaty procesu odnowy i przykład odstępów gamma}
\begin{twierdzenie}[Transformata miary odnowy]
Niech $S_0=0$ i $S_n=\sum_{k=1}^nT_k$ dla niezależnych,
ściśle dodatnich odstępów o jednakowym rozkładzie.
Zdefiniujmy lokalnie skończoną miarę na $[0,\infty)$
\[
U(A)=\sum_{n=0}^\infty\Prob(S_n\in A),\qquad
A\in\calB([0,\infty)).
\]
Jej funkcją dystrybucyjną jest $U([0,t])=1+\E N_t$.
Jeżeli $\varphi(\theta)=\E e^{-\theta T_1}$, to dla $\theta>0$
\[
\int_{[0,\infty)}e^{-\theta t}\,\mathrm dU(t)
=\frac{1}{1-\varphi(\theta)}.
\]
\end{twierdzenie}
\begin{proof}
Dla $t<\infty$ równość
$U([0,t])=1+\sum_{n\ge1}\Prob(S_n\le t)=1+\E N_t$
i poprzednie oszacowanie $\E N_t<\infty$ dowodzą
lokalnej skończoności miary. Z twierdzenia Tonellego,
nieujemności funkcji podcałkowej i niezależności odstępów
\[
\int e^{-\theta t}\,\mathrm dU(t)
=\sum_{n=0}^\infty\E e^{-\theta S_n}
=\sum_{n=0}^\infty \varphi(\theta)^n.
\]
Ponieważ $T_1>0$ p.n., mamy $0\le\varphi(\theta)<1$.
Suma szeregu geometrycznego daje tezę.
\end{proof}

\begin{zadanie}[Każde drugie zdarzenie Poissona]
Niech $P_t$ będzie procesem Poissona o intensywności $\lambda>0$.
Interesuje nas tylko co drugie zdarzenie: czasem odnowy jest
suma dwóch kolejnych odstępów wykładniczych. Obliczyć
rozkład $N_t$ i $\E N_t$.

\emph{Rozwiązanie.} Kolejne pary niezależnych
odstępów wykładniczych dają niezależne czasy
$T_n\sim\mathrm{Gamma}(2,\lambda)$ (parametry: kształt $2$,
intensywność $\lambda$). Liczba zakończonych par to
$N_t=\lfloor P_t/2\rfloor$. Jeśli $u=\lambda t$, mamy
\[
\Prob(N_t=n)=\Prob(P_t=2n)+\Prob(P_t=2n+1)
=e^{-u}\left(\frac{u^{2n}}{(2n)!}+
\frac{u^{2n+1}}{(2n+1)!}\right).
\]
Ponieważ $\lfloor k/2\rfloor=(k-\chi_{\{k\text{ nieparzyste}\}})/2$,
\[
\E N_t=\frac{\E P_t}{2}-\frac12\Prob(P_t\text{ nieparzyste}).
\]
Z funkcji tworzącej rozkładu Poissona
$\E[(-1)^{P_t}]=e^{-2u}$, skąd
$\Prob(P_t\text{ nieparzyste})=(1-e^{-2u})/2$.
Ostatecznie
\[
\E N_t=\frac{\lambda t}{2}-\frac{1-e^{-2\lambda t}}{4}.
\]
W szczególności $\E N_t/t\to\lambda/2=1/\E T_1$,
zgodnie z elementarnym twierdzeniem odnowy.
Ale $N_t$ nie jest procesem Poissona: odstępy są gamma,
a nie wykładnicze.
\end{zadanie}

\section{Podsumowanie i dalsze pytania}
W procesie odnowy niezależne są całe kolejne cykle, a nie
zwykle liczby zdarzeń w rozłącznych przedziałach. Wzór
odnowy opisuje średnią liczbę zakończonych cykli, a oba
prawa graniczne podają długookresową częstość odnowień.
Zagadnienia o wieku i czasie resztkowym wyjaśniają, dlaczego
sama liczba zdarzeń nie musi wystarczać do uzyskania własności Markowa.
Twierdzenie Blackwella oraz kluczowe twierdzenie
odnowy opisują lokalne granice miary odnowy
i jej splotów, z rozróżnieniem przypadku
arytmetycznego i niearytmetycznego.
Graniczne rozkłady wieku i czasu resztkowego
wyjaśniają paradoks inspekcji, a odpowiedni
dobór pierwszego oczekiwania daje proces
odnowy o stacjonarnych przyrostach.


% ==============================================================
% ROZDZIAŁ 10. PROCESY ROZGAŁĘZIAJĄCE
% ==============================================================
